% Problem record op_b52e28b95f7677a7. This TeX file is derived from
% database/problems_json/op_b52e28b95f7677a7.json by scripts/sync-tex.mjs; edit the JSON record
% and rerun the script rather than editing this file.
\section{Optimal output monitoring for parametric-oscillator squeezing}
\label{sec:op_b52e28b95f7677a7}

\paragraph{Problem.}
Can any causal output measurement improve the mean conditional position squeezing achieved by ideal homodyne detection?
Consider one mode with $[q,p]=i$, $a=(q+ip)/\sqrt2$, and $0\leq\chi<1/2$.
The oscillator interacts with a vacuum Markov bath at unit damping rate.
Its dynamics are Eq.~\eqref{eq:sr5-dynamics}.
\begin{equation}
H_\chi=-\frac\chi2(qp+pq),\qquad
\dot\rho=-i[H_\chi,\rho]+\mathcal D[a]\rho,\qquad
\mathcal D[a]\rho=a\rho a^\dagger-\frac12\{a^\dagger a,\rho\}.
\label{eq:sr5-dynamics}
\end{equation}
The initial oscillator state is the unconditional stationary state.
The measurement $\mathsf M$ may be adaptive or non-Gaussian and may access the output field up to the current time.
It does not apply control operations to the oscillator.
Let $\rho_c(t)$ be the oscillator state conditioned on the measurement record.
The mean is over all records, without postselection.
Define the optimum by Eq.~\eqref{eq:sr5-cost}.
\begin{equation}
s_{\mathrm{all}}(\chi)=\inf_{\mathsf M}\liminf_{t\to\infty}\mathbb E_{\mathsf M}\!\left[2\operatorname{Var}_{\rho_c(t)}q\right].
\label{eq:sr5-cost}
\end{equation}
Does $s_{\mathrm{all}}(\chi)=1-2\chi$ hold throughout the stated range?

\subsection*{Status}

Solved

\subsection*{Source}

This is a formulation for a fixed-quadrature mean cost in the parametric oscillator of Genoni, Lami, and Serafini, Sec.~6.1 \sourcecite{ref:sr5-genoni}{GLS16}. The extension of its monitoring optimization to arbitrary non-Gaussian output measurements is editorial.

\subsection*{Progress}

\begin{itemize}
  \item At unit damping and zero temperature, Sec.~6.1, Eqs.~(77)--(78) of the arXiv PDF, give the stationary unconditional and ideal-homodyne covariances in Eq.~\eqref{eq:sr5-covariances} \sourcecite{ref:sr5-genoni}{GLS16}. Choose the homodyne phase whose output signal is $a+a^\dagger=\sqrt2q$. The covariance convention is $V_{jk}=\langle\{R_j-\langle R_j\rangle,R_k-\langle R_k\rangle\}\rangle$ for $R=(q,p)^{\mathsf T}$.
\begin{equation}
\begin{aligned}
 V_{\mathrm{unc}}&=\operatorname{diag}\!\left(\frac1{1+2\chi},\frac1{1-2\chi}\right),\\
 V_{\mathrm{hom}}&=\operatorname{diag}\!\left(1-2\chi,\frac1{1-2\chi}\right).
\end{aligned}
\label{eq:sr5-covariances}
\end{equation}

  \item A direct ensemble argument proves optimality even for non-Gaussian monitoring. Write $V_{q,c}=2\operatorname{Var}_{\rho_c}q$ and $V_{p,c}=2\operatorname{Var}_{\rho_c}p$. Robertson uncertainty gives $V_{q,c}V_{p,c}\geq1$. Jensen's inequality and the law of total variance then imply Eq.~\eqref{eq:sr5-lower} at every time.
\begin{equation}
\mathbb E[V_{q,c}]\geq\mathbb E[1/V_{p,c}]\geq\frac1{\mathbb E[V_{p,c}]}\geq\frac1{(V_{\mathrm{unc}})_{pp}}=1-2\chi.
\label{eq:sr5-lower}
\end{equation}
The unconditional state stays stationary because measurements act only on the output. Combining Eq.~\eqref{eq:sr5-lower} with the homodyne value in Eq.~\eqref{eq:sr5-covariances} proves $s_{\mathrm{all}}(\chi)=1-2\chi$. This last argument is supplied here; it is not attributed to the paper.
\end{itemize}

\subsection*{References}

\begin{enumerate}[label={},leftmargin=4.8em,labelsep=0.5em]
  \footnotesize
  \item[\textup{[GLS16]}]\label{ref:sr5-genoni}
  M. G. Genoni, L. Lami, and A. Serafini, ``Conditional and Unconditional Gaussian Quantum Dynamics,'' \emph{Contemporary Physics} \textbf{57}, 331--349 (2016). \href{https://doi.org/10.1080/00107514.2015.1125624}{doi:10.1080/00107514.2015.1125624}; \href{https://arxiv.org/abs/1607.02619}{arXiv:1607.02619}.
\end{enumerate}

\subsection*{Comment}

The fixed-quadrature mean cost is completely resolved by the published homodyne solution and the direct lower-bound argument above. The latter is not a separately peer-reviewed theorem. Outcome-dependent quadrature choices, postselected costs, and feedback control on the oscillator define different optimization problems.

\subsection*{Field}

Quantum Resource Theory

\subsection*{Topic}

Gaussian quantum information

\subsection*{ID}

\texttt{op\_b52e28b95f7677a7}
